\documentclass[onecolumn]{article}

\usepackage{tikz}
\usetikzlibrary{arrows}

\usepackage{graphicx}
\usepackage{verbatim}
\usepackage{subfigure}
\usepackage{multicol}
\usepackage{hyperref}
\usepackage{framed}
\usepackage[T1]{fontenc}
\usepackage[utf8] {inputenc} %tildes

\usepackage[all]{xy}
\usepackage{amsmath}
\usepackage{amssymb}
\usepackage{amsfonts}
\usepackage{lscape}
\usepackage{amsthm}
\usepackage[mathscr]{euscript}
\usepackage[normalem]{ulem}  %% tachar texto

\usepackage{fancyhdr}
\usepackage{multirow}
%\usepackage{booktabs}

%\usepackage{tikz}


\usepackage{layout}
\setlength{\hoffset} {0cm}        % 1
\setlength{\voffset} {0cm}        % 2
\setlength{\oddsidemargin}{0cm}   % 3
\setlength{\evensidemargin}{0cm}  % 4
\setlength{\topmargin}{1cm}       % 5
\setlength{\headheight}{4cm}      % 6
\setlength{\headsep}{1cm}         % 7
\setlength{\textwidth}{17cm}      % 8
\setlength{\textheight}{17cm}     % 9
\setlength{\marginparsep}{0cm}    % 10
\setlength{\marginparwidth}{0cm}  % 11
\setlength{\footskip}{0cm}        % 12


%%% Definiciones para su Gu�a de laboratorio%%%
\newcommand{\profesor}{Leonardo Rodríguez Mújica}
\newcommand{\correo}{lrodriguez@mail.unicundi.edu.co}
\newcommand{\nucleotematico}{Antenas}
\newcommand{\codigo}{622007741}
\newcommand{\monitor}{}
\newcommand{\cormon}{}

\newcommand{\marco}[2]
{
{\fboxsep 10pt
    \begin{framed}

\section{#1}
#2
\end{framed}
}
}

\newcommand{\continua}[1] { 

\begin{center}
{\fboxsep 10pt
\fcolorbox{black}{white}{
\begin{minipage}[t]{16cm}
#1
\end{minipage}
}}
\end{center}
}

\newcommand{\objetivos}[1] { \marco{Objetivos}{#1} }
\newcommand{\materiales}[1] { \marco{Materiales sugeridos para construccion antena Yagi}{#1} }
\newcommand{\planteamiento}[1] { \marco{Planteamientos y procedimiento}{#1} }
\newcommand{\metodologia}[1] { \marco{Metodología y evaluación}{#1} }
\newcommand{\recomendaciones}[1] { \marco{Recomendaciones}{#1} }
\newcommand{\bibliografia}[1] { \marco{Bibliografia}{#1} }

%%% Digitado por el usuario {\tt }
%%% Click en alguna parte {\sl }
%%% ingl�s {\em }

\cfoot{
Diagonal 18 No 20-29  Fusagasugá- Cundinamarca Teléfonos: 8732512/35 - Fax \\
Página Web: \href{www.unicundi.edu.co}{
www.unicundi.edu.co} - Correo Electrónico \href{mailto:unicundi@unicundi.edu.co}{unicundi@unicundi.edu.co}
}

\fancyhead{
\begin{tabular}{|c|c|c|c|}
\hline
  % after \\: \hline or \cline{col1-col2} \cline{col3-col4} ...
  \hspace{3cm}  & \hspace{8cm}                        & \hspace{2cm}    & \hspace{2cm}     \\
                & \textbf{UNIVERSIDAD DE CUNDINAMARCA}  & FECHA DE          &  CÓDIGO            \\
                & \textbf{FACULTAD DE INGENIERÍA}       & EMISIÓN           &                    \\
                & \textbf{INGENIERÍA ELECTRÓNICA}       &                   &                    \\
                &                                       &  \date{\today}    &                    \\ \cline{2-4}
                &                                       &                   &                    \\
                & \textbf{GUÍA DE LABORATORIO  }        & VERSION No.1.1    &  PÁGINA \thepage   \\
                &                                       &                   &                    \\
  \hline
\end{tabular}
}

\fancyhead[L]{
\hspace{0.4cm}
\vspace{0.4cm}
\includegraphics[scale=0.12]{logo.eps}
}

\title{\textbf{FORMATO GUÍA DE LABORATORIO \\ UNIVERSIDAD DE CUNDINAMARCA \\ FACULTAD DE INGENIERÍA \\ PROGRAMA INGENIERIA ELECTRÓNICA}}
\author{\textbf{Jose David Alvarado Moreno} \\ \\ \huge{\textbf{\LaTeX}}}
\date{\today}

\begin{document}

\pagestyle{fancy}

%%\maketitle

\begin{center}
\begin{tabular}{|l l|l l|}
\hline
\multirow{2}{*}{Asignatura:} &\multirow{2}{*}{\textbf{\nucleotematico}}  & \multirow{2}{*}{Código:}  & \multirow{2}{*}{\codigo}  \\
                             &                                         &                          &                    \\ \hline
\multirow{2}{*}{Docente:}    &\multirow{2}{*}{\textbf{\profesor}}   & \multirow{2}{*}{E-mail:} & \multirow{2}{*}{\href{mailto:\correo}{\correo}}  \\
                             &                                         &                          &                     \\ \hline
\multirow{2}{*}{Monitor:}    &\multirow{2}{*}{\textbf{\monitor}} & \multirow{2}{*}{E-mail:} & \multirow{2}{*}{\href{mailto:\cormon}{\cormon}}   \\
                             &\hspace{6cm}                             &                          & \hspace{6cm}        \\ \hline
\end{tabular}
\end{center}

\begin{center}
\section*{PRACTICA No 2}
\addcontentsline{toc}{section}{PRACTICA No 2}

\section*{Diseño e implemetacion antena YAGI}

\addcontentsline{toc}{section}{construccion antena}
\end{center}
%%%--- Colocar objetivos acá ----%%%%

\objetivos{

\begin{itemize}
\item  Diseñar antena Yagi para una emisora universitaria de la
    ciudad de Bogotá utilizando los software (YagiModel y XNEC2C)
\item  Implementar y probar la antena Yagi
\item Analizar los  diseños con los dos programas (YagiModel y xnec2c)
\end{itemize}
}

%%%----- Listar los materiales acá ------%%%
\materiales{

\begin{enumerate}
  \item Tubos de aluminio
  \item Cable coaxial
  \item XNEC2c y YagiModeler    
\end{enumerate}
}


%%% ------- colocar planteamiento y procedimiento

%% 1. simular la yagi con yagi model
%% 2. con las fotos mostrar como debe quedar
%% 3. comparar con la simulación de xnec2c


%% 1.  hacer el balum
%% averiguar bien el tipo de cable con que se va a hacer
%% 2. probar la antena con el equipo
%% 3. enseñar a conectarse en linux
%% 4. ver la intensidad de la señal con wavemon


%% grupo especial para práctica de patron de radiación. Replicar la práctica de méxico.
%% 1. probar dipolo en equipo
%% 2. probar galvanometro con diferentes access points y equipos de forma incremental.

\planteamiento{

\begin{enumerate}

%% Diseño antena emisora universitaria

    \item Partiendo de la emisora universitaria que le correspondió y del
        dipolo diseñado en la práctica pasada diseñe una antena yagi cuyo
        objetivo es captar la señal de la emisora de una forma adecuada.
        para ello siga esta metodología \begin{enumerate}       
%% simulación modeler
        \item Ejecute el programa {\tt yagui-modeler} descomprimiendo el
            archivo {\tt yagi-modeler.tar.gz} que se encuentra en la
            plataforma; abra {\tt index.html} y presione el botón de {\em
            Click Here for the Yagi Modeler}. Si tiene problemas con la
            seguiridad de Java, cambie esta configuración.

%% simulación modeler
        \item  En el menú {\tt examples} escoja un diseño del tipo {\tt
            K1FO 70 cm}, procurando quecada grupo implemente una antena
            diferente. 
            adapte este diseño para que trabaje a la frecuencia portadora
            de su emisora siguiendo los pasos descritos en "{\em This
            applet can be used to scale the example yagis or any others to
        new frequencies}. 

%% simulación modeler
        \item Seleccionando el ejemplo deseado trabajar sobre él, en éste
            recuadro nos permite ajustar la frecuencia  a trabajar y el
            diámetro del tubo a utilizar. También pueden observar
            impedancia y ganacias.

%% simulacion
        \item Documente: ganancia, impedancia, patrón de radiación de la
            antena.  También la longitud y posición de cada elemento.

    \end{enumerate}
    y consulte lo siguientes parámetros de su emisora para hacer un
    estimativo de la potencia recibida por nuestra antena.

    \begin{enumerate}
        \item Ubicación geográfica de los transmisores
        \item potencia de transmisión 
        \item Escoja una posible antena que utilice el transmisor (consulte
            antenas típicas de un transmisor de FM)
        \item Con la localización geográfica de Fusagasugá trate de
            consultar o calcular
            el ángulo al cual habría que apuntar la antena para una
            recepción óptima.
    \end{enumerate}
%% simulación Xnec2c
   
\item Dibuje y simule esta misma antena con {\tt xnec2c}. Comparar los
    diseños de las antenas basándose en los resultados de la simulación
    tener en cuenta patrón de radiación, potencias, ganancias, impedancias
    etc. Describa fortalezas y debilidades de cada una de las
    herramientas.



%% implementar
\item  Implementar  antena Yagi. Tenga en cuenta la lista de materiales
    suministrada por el programa con los tamaños distancias y material.
    Seleccione un buen soporte (no conductor) para los elementos de la
    Yagi. Haga pruebas de su emisora, evalúe el ancho de banda de la antena
    sintonizando emisoras en todo el dial.

%% implementar
\item Según la implementación hecha y las simulaciones aclare las
    siguientes cuestiones: cuál es el ancho de banda de la antena; según el
    cable utilizado cuál será el coeficiente de reflexión; comparelo con lo
    obtenido en xnec2c si no coinciden corrija la simulación;
    si existiera
    línea de vista entre el transmisor y el receptor cuál sería la potencia
    y el voltaje recibidos.


%Los parámetros de diseño empleados para esta antena son: frecuencia 2.4GHz, dipolo $\lambda/2$, reflectores y directores (seleccionarlos a su criterio)
%\item  Teniendo conocimiento e interactuado con el simulador XNEC2C cargar ejemplo de la antena Yagi, y desde esta este diseñar su modelo de antena
%\item Utilizando el software YagiModel diseñar su modelo de antena tenga en cuenta los siguientes pasos

%-Cargue un ejemplo de la antena como se ilustra en la imagen
%%cargar imagen ya.1jpg


%%cargar imagen ya.2jpg

%-Seleccionamos la conductividad del cable
%%cargar imagen ya5.jgp

%-Clikcleamos el boton calcular y observaremos el patron de radiacion y diferentes medidas que este nos arroja
%%cargar imagen ay.3jpg

%-Buscamos en file y podemos obtener la lista de los elementos la posicion y longitud a los cuales se deben ubicar
%%cargar imagen ya.4jpg
%% funcionamiento simulación
%\item Obtener patron de radiacion de las diferentes herramientas (YagiModel XNEC2C)


%\item Diseñar e implementar el balun para adaptar impedancias
%% xnec2c, diseño
%\item  Implementar  antena Yagi 

%%


\end{enumerate}
}

%%% --- Colocar acá la metodología y evaluación
\metodologia{

Presentar un informe con los resultados obtenidos en la práctica que se
evaluará  teniendo en cuenta la siguiente rúbrica.

\newcommand{\med}{0.22}
\begin{tabular}{|p{\med\textwidth}|p{\med\textwidth}|p{\med\textwidth}|p{\med\textwidth}|p{\med\textwidth}|}
\hline
criterio&Excelente(90\%)&Bueno(70\%)&Deficiente(10\%)\\
\hline
Interaccion con XNEC2C y YagiModel(25\%)&Documentación de la simulación antena&Genera documento sin las principales tareas&No documenta actividades con XNEC2C\\
\hline
Funcionamiento de la antena Yagi(25\%)&Realiza todas las tareas con el simulador&Realiza sólo algunas simulaciones&No emplea el simulador\\
\hline
Diseño e implementación de la antena(25\%)&Implementa la antena y la prueba adecuadamente&Implmentación y pruebas deficientes&No realiza las actividades\\
\hline
Forma:presentación, redacción ortografía (10\%)&Sin errores y bien
presentado&Presenta deficiencias&No elabora documento\\
\hline
Conclusiones(15\%)&claras, concisas y relevantes&Falta especificar resultados&Irrelevantes e intrasendentes\\
\hline
\end{tabular}
}
%%% --- colocar acá las recomendaciones ---%%%
\recomendaciones{

\begin{itemize}
\item Leer sobre el software YagiModel y diseño de antenas Yagi para realizar la práctica
\item Seleccionar calibres de cable adecuado tanto para simulación como para la implementación

\end{itemize}
}

%%% ---- listar las referencia bibliogr�ficas
\bibliografia{

\begin{itemize}\itemsep=-4pt
  \item Meserve, Martin E. VHF/UHF Yagi Antenna Feed Design.\url{http://www.k7mem.com/Electronic_Notebook/antennas/yagi_vhf_feed.html}
  \item Xnec2c User Manual.\url{http://www.qsl.net/5b4az/pkg/nec2/xnec2c/doc/xnec2c.html }
  \item Antena Yagi
Networks.\url{http://arieldx.tripod.com/estaciondx/proyectos/yagifm.htm}
\item Schmidt, Kevin. Software YagiModel.\url {http://fermi.la.asu.edu/w9cf/yagipub/}
\item \url{http://tf.nist.gov/general/pdf/451.pdf} \url{http://www.researchgate.net/post/Regarding_Yagi-Uda_antenna}
 
\end{itemize}
}


\end{document} 
